\section{Game Theory dan Parameter Bias Lokal $h_i$}

\subsection{Konstruksi Matriks Payoff Historis}
Matriks imbalan (\textit{payoff matrix}) $V_i$ disusun berdasarkan rata-rata imbal hasil historis aset $i$ pada setiap kemungkinan konfigurasi status dengan aset $j$. Elemen-elemen $R_{xy}$ didefinisikan sebagai:
\begin{align}
R_{UU} &= \frac{1}{n_{UU}} \sum_{t \in T_{UU}} r_{i,t} \\
R_{UD} &= \frac{1}{n_{UD}} \sum_{t \in T_{UD}} r_{i,t} \\
R_{DU} &= \frac{1}{n_{DU}} \sum_{t \in T_{DU}} r_{i,t} \\
R_{DD} &= \frac{1}{n_{DD}} \sum_{t \in T_{DD}} r_{i,t}
\end{align}
di mana $T_{xy}$ adalah himpunan waktu $t$ yang memenuhi kondisi $(s_{i,t}, s_{j,t}) \in \{U, D\}^2$. Matriks $V_i$ kemudian dinyatakan sebagai:
\begin{equation}
V_i = \begin{pmatrix} R_{UU} & R_{UD} \\ R_{DU} & R_{DD} \end{pmatrix}
\end{equation}

\subsection{Ekspansi Fungsi Utilitas Biner}
Pemetaan dari matriks diskret ke fungsi utilitas kontinu $V_i(s_i, s_j)$ menggunakan variabel spin $s \in \{+1, -1\}$ dilakukan melalui ekspansi biner:
\begin{equation}
\begin{split}
V_i(s_i, s_j) =& \frac{1+s_i}{2} \left[ \frac{1+s_j}{2} R_{UU} + \frac{1-s_j}{2} R_{UD} \right] \\
&+ \frac{1-s_i}{2} \left[ \frac{1+s_j}{2} R_{DU} + \frac{1-s_j}{2} R_{DD} \right]
\end{split}
\label{eq:v_function_expansion_detail}
\end{equation}

\subsection{Normalisasi Probabilitas dan Distribusi Strategi}
Frekuensi kejadian historis dipetakan menjadi distribusi probabilitas untuk menentukan strategi agen. Probabilitas gabungan $P(s_i, s_j)$ dihitung dengan membagi frekuensi kemunculan konfigurasi $n(s_i, s_j)$ terhadap total agregat observasi $N$:
\begin{equation}
P(s_i, s_j) = \frac{n(s_i, s_j)}{N}
\end{equation}
Probabilitas marginal untuk strategi agen $j$, yang merepresentasikan distribusi preferensi strateginya, diturunkan melalui penjumlahan parsial:
\begin{equation}
P(s_j) = \sum_{s_i \in \{+1, -1\}} P(s_i, s_j)
\end{equation}
Sifat normalisasi menjamin bahwa total probabilitas gabungan bernilai satu ($\sum P(s_i, s_j) = 1$).

\subsection{Ekspektasi Payoff dan Strategi Nash}
Nilai harapan imbalan (\textit{expected payoff}) dihitung dengan merata-ratakan utilitas terhadap distribusi probabilitas strategi lawan $P(s_j)$:
\begin{equation}
E[i|s_i] = \sum_{s_j \in \{+1, -1\}} P(s_j) \cdot V(s_i, s_j)
\end{equation}
Untuk status spesifik $s_i = +1$ (Naik) dan $s_i = -1$ (Turun), diperoleh marginal payoff:
\begin{align}
E[i|+] &= P(+1)V(+1,+1) + P(-1)V(+1,-1) \\
E[i|-] &= P(+1)V(-1,+1) + P(-1)V(-1,-1)
\end{align}

\subsection{Ekuivalensi ke Parameter Bias Lokal $h_i$}
Berdasarkan aksioma pemetaan energi terhadap utilitas $H \equiv -E$, Hamiltonian lokal $H_{lokal}(s_i) = -h_i s_i$ harus memenuhi kondisi kesetimbangan:
\begin{align}
H_{lokal}(+1) &= -h_i \equiv -E[i|+] \\
H_{lokal}(-1) &= +h_i \equiv -E[i|-]
\end{align}
Melalui evaluasi selisih energi sistem:
\begin{equation}
\Delta H = H_{lokal}(-1) - H_{lokal}(+1) = (+h_i) - (-h_i) = 2h_i
\end{equation}
Substitusi nilai utilitas marginal menghasilkan:
\begin{equation}
2h_i = (-E[i|-]) - (-E[i|+]) = E[i|+] - E[i|-]
\end{equation}
Sehingga formulasi akhir parameter bias lokal $h_i$ adalah:
\begin{equation}
h_i = \frac{E[i|+] - E[i|-]}{2}
\end{equation}
Nilai ini merepresentasikan kecenderungan intrinsik aset untuk berada pada status tertentu berdasarkan interaksi strategis historisnya.
